\chapter{Adaptive stationäre Wärme}\label{api:adaptive}
\index{Adaptivität}\index{Kantenstern-Bisektion}
Der additive M10-Vertrag gilt ausschließlich für stationäre Wärme auf affinen,
konformen Tet4-Netzen. \path{HeatProblem} enthält ein Material oder eine geordnete
Materialfolge je Zelle, Dirichlet-Werte je Facettenzeile und optional Quelle,
Auswärtsflüsse und Robin-Daten. SI-Einheiten und Vorzeichen entsprechen dem
stationären Wärmesolver. Kinder erben ihre Material-/Randzuordnung exakt;
räumliche Callables werden an den neuen physikalischen Punkten ausgewertet.
\path{AdaptiveOptions} validiert nichtnegative endliche Indikatortoleranz,
positive ganze Schritt-/DOF-Budgets, Markierungsanteil und Stagnationskonfiguration.
\begin{lstlisting}
run = solve_adaptive_heat(mesh, problem, AdaptiveOptions(
    tolerance=0.1, max_steps=8, max_dofs=20000, fraction=0.5))
run.export("results/adaptive_heat.json")
\end{lstlisting}
Endstatus ist \path{tolerance}, \path{dof_budget}, \path{step_budget},
\path{stagnation} oder \path{error}. Nach einem Solve gilt die Priorität Toleranz,
Schrittbudget, Stagnation, dann Verfeinerungs-/DOF-Budget. Fehler bewahren alle
bereits abgeschlossenen Stufen. Eine Indikatortoleranz ist keine garantierte
Schranke des tatsächlichen Energiefehlers. Ungültige Optionen werfen
\path{ValidationError}; unerwartete Softwarefehler dürfen nicht als Erfolg gelten.

\path{AdaptiveStage} speichert Mesh, Rohresultat, Indikator, Markierung und
Elternzeilen. Unveränderliche Arrays und kopierte Mappings verhindern nachträgliche
Änderungen der Evidenz. JSON-Schema 1 enthält Optionen, Status, Netze samt
Tags/Gruppen, SHA-256, DOFs, Felder, Reaktionen, Residuen, Bilanz und Provenienz.
Zur erneuten Ausführung von Callables bleibt das Problemskript erforderlich.
Export überschreibt den ausdrücklich angegebenen Pfad; eine Resume-Funktion
oder TOML-Adaptivitätskonfiguration ist nicht zugesagt.

\path{refine_tet4} teilt die längste Eingabekante jeder markierten Zelle,
sortiert nach Knotenpaaren. Jeder vollständige Kantenstern und seine Randdreiecke
teilen denselben Mittelpunkt; unmarkierte Nachbarn können mitverfeinert werden.
\path{Refinement} liefert Kind-Eltern-Zeilen und verwirft externe IDs.
Die Gruppen-/Geometrievalidatoren sind Forschungsbausteine, keine allgemeinen
Schnittdetektoren. Keine formale gleichmäßige Formregularitätsgarantie,
Koarsening-, Glättungs-, Tet10-, gekoppelte oder transiente Adaptivität.
